\section{Fibonacci Number}
\label{sec:dp-fibonacci}

\problemstatement{Compute the $n$-th Fibonacci number, defined by
$F(0)=0$, $F(1)=1$, and $F(n) = F(n-1) + F(n-2)$ for $n \ge 2$.}

\begin{patternbox}
\emph{``Each term is defined by combining a small, fixed number of
earlier terms''} is the simplest possible 1D DP keyword pattern, and
Fibonacci is its archetype: the recurrence $F(n) = F(n-1)+F(n-2)$ is
already explicitly given, so no derivation is needed -- only the
discipline of carrying it through all four stages described in
Section~\ref{sec:dp-intro}.
\end{patternbox}

\subsection*{Approach 1: Recursion}

The recurrence translates directly into a function with two base cases
($n=0$ and $n=1$) and one recursive case.

\begin{lstlisting}
#include <bits/stdc++.h>
using namespace std;

int fibRecursive(int n) {
    if (n <= 1) return n;
    return fibRecursive(n - 1) + fibRecursive(n - 2);
}

int main() {
    cout << fibRecursive(10) << endl; // 55
    return 0;
}
\end{lstlisting}

\begin{complexitybox}[Approach 1 Complexity]
\textbf{Time.} Each call spawns two more, and the recursion tree has
depth $n$, giving $O(2^n)$ -- more precisely $O(\phi^n)$ where $\phi$ is
the golden ratio, but exponential either way.
\textbf{Space.} $O(n)$ for the recursion stack depth.
\end{complexitybox}

\subsection*{Approach 2: Memoization (Top-Down DP)}

\begin{ideabox}[Why $\textsc{Fib}(k)$ Is Computed Repeatedly]
Expanding $\textsc{Fib}(5) = \textsc{Fib}(4) + \textsc{Fib}(3)$: the call
to $\textsc{Fib}(4)$ itself expands to $\textsc{Fib}(3)+\textsc{Fib}(2)$,
so $\textsc{Fib}(3)$ is computed once directly and once more inside
$\textsc{Fib}(4)$'s expansion -- and this doubling compounds at every
level. Caching each $\textsc{Fib}(k)$'s result the first time it is
computed, in an array $\textit{memo}$ indexed by $k$, means every later
request for the same $k$ is an $O(1)$ lookup instead of a full
re-expansion.
\end{ideabox}

\begin{lstlisting}
#include <bits/stdc++.h>
using namespace std;

int fibMemo(int n, vector<int>& memo) {
    if (n <= 1) return n;
    if (memo[n] != -1) return memo[n];

    memo[n] = fibMemo(n - 1, memo) + fibMemo(n - 2, memo);
    return memo[n];
}

int main() {
    int n = 10;
    vector<int> memo(n + 1, -1);
    cout << fibMemo(n, memo) << endl; // 55
    return 0;
}
\end{lstlisting}

\begin{complexitybox}[Approach 2 Complexity]
\textbf{Time.} Each of the $n+1$ distinct sub-problems
$\textsc{Fib}(0),\ldots,\textsc{Fib}(n)$ is computed exactly once,
each doing $O(1)$ work beyond its (now-cached) recursive calls, giving
$O(n)$.
\textbf{Space.} $O(n)$ for the memo array, plus $O(n)$ for the
recursion stack -- still $O(n)$ overall, but for a different reason than
Approach 1 (here it's the cache, not the redundant branching, that
dominates the asymptotic class, even though the stack depth alone was
already $O(n)$ in Approach 1 too).
\end{complexitybox}

\subsection*{Approach 3: Tabulation (Bottom-Up DP)}

\begin{ideabox}[Filling the Table in the Opposite Direction]
Memoization computes $\textsc{Fib}(n)$ by recursing \emph{downward}
toward the base cases, caching as it unwinds. Tabulation instead starts
\emph{at} the base cases and fills an array $\textit{dp}$ \emph{upward}
toward $n$ in an explicit loop -- computing the exact same values, in the
same dependency order, but without any function call overhead.
\end{ideabox}

\begin{lstlisting}
#include <bits/stdc++.h>
using namespace std;

int fibTabulation(int n) {
    if (n == 0) return 0;

    vector<int> dp(n + 1);
    dp[0] = 0;
    dp[1] = 1;
    for (int i = 2; i <= n; i++) {
        dp[i] = dp[i - 1] + dp[i - 2];
    }
    return dp[n];
}

int main() {
    cout << fibTabulation(10) << endl; // 55
    return 0;
}
\end{lstlisting}

\subsection*{Dry run of the tabulation table}

\dryrun{Take $n=6$.}

\begin{center}
\begin{tabular}{c|c|c|c|c|c|c|c}
$i$ & 0 & 1 & 2 & 3 & 4 & 5 & 6 \\ \hline
$\textit{dp}[i]$ & 0 & 1 & 1 & 2 & 3 & 5 & 8
\end{tabular}
\end{center}

Each entry from $i=2$ onward is filled as $\textit{dp}[i-1]+\textit{dp}[i-2]$:
e.g.\ $\textit{dp}[4] = \textit{dp}[3]+\textit{dp}[2] = 2+1 = 3$. The
answer, $\textit{dp}[6]=8$, matches $F(6)=8$.

\begin{complexitybox}[Approach 3 Complexity]
\textbf{Time.} $O(n)$ -- one pass filling $n-1$ table entries.
\textbf{Space.} $O(n)$ for the $\textit{dp}$ array; \textbf{no}
recursion stack at all, unlike Approaches 1 and 2.
\end{complexitybox}

\subsection*{Approach 4: Space Optimization}

\begin{ideabox}[Only the Last Two Entries Are Ever Needed]
Computing $\textit{dp}[i]$ only reads $\textit{dp}[i-1]$ and
$\textit{dp}[i-2]$ -- every entry before $i-2$ is permanently dead the
moment $\textit{dp}[i]$ is computed. Replacing the full array with two
rolling scalars, updated each iteration, preserves correctness while
cutting space from $O(n)$ to $O(1)$.
\end{ideabox}

\begin{lstlisting}
#include <bits/stdc++.h>
using namespace std;

int fibSpaceOptimized(int n) {
    if (n == 0) return 0;

    int prev2 = 0, prev1 = 1;
    for (int i = 2; i <= n; i++) {
        int current = prev1 + prev2;
        prev2 = prev1;
        prev1 = current;
    }
    return prev1;
}

int main() {
    cout << fibSpaceOptimized(10) << endl; // 55
    return 0;
}
\end{lstlisting}

\begin{complexitybox}[Approach 4 Complexity]
\textbf{Time.} $O(n)$, unchanged from tabulation.
\textbf{Space.} $O(1)$ -- only two scalar variables, regardless of how
large $n$ is.
\end{complexitybox}

\begin{takeawaybox}
Every later 1D DP problem in this chapter follows exactly this
progression. The only thing that will change from here on is the
\emph{recurrence itself} -- how many earlier states each step depends on,
and how they are combined (summed, for counting; minimized or maximized,
for optimization). The four-stage scaffolding shown here is reused
verbatim.
\end{takeawaybox}
